\section{总结与延伸}

\subsection{讲者的核心总结}

Jiajun Wu 在课程最后回顾了 3D 视觉研究的发展轨迹：

\begin{enumerate}
    \item \textbf{表示方式决定了方法论}：从体素到点云到隐式函数，每种表示催生了不同的深度学习方法
    \item \textbf{显式与隐式的互补}：采样容易 vs 查询容易，二者各有所长
    \item \textbf{2D-3D 联合学习是趋势}：3D 数据稀缺，但 2D 数据丰富。通过可微渲染和隐式表示，可以有效利用 2D 先验
    \item \textbf{NeRF 和 3DGS 改变了游戏规则}：从 2D 图像直接学习 3D 场景，不再需要显式的 3D 监督
\end{enumerate}

\subsection{全课知识图谱}

\begin{center}
\begin{tikzpicture}[node distance=0.8cm, every node/.style={draw, rounded corners, minimum width=3.5cm, minimum height=0.7cm, font=\small}]
\node (repr) {3D 表示方式};
\node (exp) [below left=1cm and 0.5cm of repr] {显式表示};
\node (imp) [below right=1cm and 0.5cm of repr] {隐式表示};
\node (dl) [below=2.5cm of repr] {3D 深度学习方法};
\node (nerf) [below=of dl] {NeRF / 3DGS};
\node (gen) [below=of nerf] {3D 生成模型};
\draw[->, thick] (repr) -- (exp);
\draw[->, thick] (repr) -- (imp);
\draw[->, thick] (exp) -- (dl);
\draw[->, thick] (imp) -- (dl);
\draw[->, thick] (dl) -- (nerf);
\draw[->, thick] (nerf) -- (gen);
\node[draw=none, left=0.5cm of exp, text width=4cm, font=\scriptsize] {点云、网格、参数曲面\\采样容易，查询困难};
\node[draw=none, right=0.5cm of imp, text width=4cm, font=\scriptsize] {SDF、占用函数、水平集、体素\\查询容易，采样困难};
\end{tikzpicture}
\end{center}

\subsection{关键 Takeaways}

\begin{importantbox}{五条核心原则}
\begin{enumerate}
    \item \textbf{没有统一的 3D ``像素''}：不同表示各有优劣，选择取决于任务
    \item \textbf{显式 vs 隐式是核心对偶}：采样容易与查询容易的权衡贯穿始终
    \item \textbf{PointNet 的简洁之道}：排列不变 + 对称聚合，简单而强大
    \item \textbf{NeRF 用 2D 图像学 3D}：通过可微体渲染，无需 3D 监督
    \item \textbf{数据差距驱动方法创新}：3D 数据远少于 2D，利用 2D 先验是关键策略
\end{enumerate}
\end{importantbox}

\subsection{拓展阅读}

\begin{itemize}
    \item PointNet: \url{https://arxiv.org/abs/1612.00593}
    \item DeepSDF: \url{https://arxiv.org/abs/1901.05103}
    \item Occupancy Networks: \url{https://arxiv.org/abs/1812.03828}
    \item NeRF: \url{https://arxiv.org/abs/2003.08934}
    \item 3D Gaussian Splatting: \url{https://arxiv.org/abs/2308.14737}
    \item DreamFusion: \url{https://arxiv.org/abs/2209.14988}
    \item ShapeNet: \url{https://shapenet.org/}
    \item Objaverse: \url{https://objaverse.allenai.org/}
\end{itemize}

\end{document}
